% ===================================================================
% DISCUSSION AND LIMITATIONS. ~0.5–1 page. Four honest bullets.
% Written from known open items in notes/open_questions.md and
% target.md §4. No overclaiming.
% ===================================================================

\paragraph{Cross-entropy extension.} Theorem~\ref{thm:general}
is proved in the quadratic-loss setting with $H_t$ a
positive-semidefinite rank-$r$ matrix. For LLM cross-entropy
loss, the natural $H_t$ is the Fisher information matrix, which
is rank-$r$ PSD near a LoRA fine-tuned point but not a projector.
Our proof uses only the per-task identity $H_tH_t^+H_t = H_t$
and the $O(d)$-equivariance of the hard distribution, both of
which survive verbatim under the Fisher-metric interpretation
(the $D_t$-freeness of Step 5 is the key). A fully rigorous
MSE-to-CE bridge via local quadratic approximation requires
controlling the second-order remainder across the $H_t$-ball, and
we can be precise about what that control needs. Writing the
task-$t$ CE loss as $L_t(\theta_0 + w) = L_t(\theta_0 + \tau_t) +
\tfrac12 (w-\tau_t)^\top H_t (w-\tau_t) + R_t(w)$ with $H_t$ the
Fisher at $\tau_t$, the Taylor remainder satisfies
$|R_t(w)| \leq \tfrac16 M_t \|w-\tau_t\|^3$, where $M_t$ bounds the
operator norm of the third derivative of $L_t$ over the segment.
Two facts make this remainder controllable in our regime. First,
the rate budget forces $\|w - \tau_t\|$ small: at the rates we
study the merged delta lies within an $O(B\,2^{-R/\deff})$ ball of
$\tau_t$ on the active subspace, so $|R_t| = O(M_t B^3
2^{-3R/\deff})$ is third-order in the same small quantity the
quadratic term is second-order in, and is dominated once $R$ is
above a model-dependent constant. Second, $M_t$ is finite for the
softmax CE loss with bounded logits (the third derivative of
log-sum-exp is uniformly bounded), so the bound is not vacuous.
This third derivative does not grow with the ambient dimension
$d$: for a single attention head with logits bounded by $L$, the
softmax-CE third derivative decomposes per token and per output
coordinate, giving $M_t \leq \mathrm{poly}(L)$ independent of $d$,
so the remainder $O(M_t B^3 2^{-3R/\deff})$ stays genuinely
third-order even at the $d \sim 4096$ scales of interest. We state
this as an informal estimate; a fully quantitative theorem -- making $M_t$ explicit for a
transformer LoRA and propagating it through the rate-distortion
argument -- is the natural next step and is what would let the
bound be quoted directly against measured CE-NLL; we give the
remainder-bounded sketch in the supplementary material and leave
the explicit-$M_t$ version for future work. The empirical section
sidesteps the issue by reporting the CE excess
$\Delta L_t$ directly rather than inferring it from the quadratic
surrogate.

\paragraph{The shared-$V$ exponent gap.} In the floor-positive
regime (e.g.\ shared or substantially-overlapping $V_t$), the
explicit linear construction of \S\ref{sec:achv} achieves an
excess distortion decay of $2^{-2R/(\deff+|A|-1)}$, slower than
the lower bound's $2^{-2R/\deff}$. We verified numerically that
a random-codebook (non-linear) encoder achieves slopes strictly
steeper than the linear construction, so the truth lies
\emph{strictly between} the two bounds. Closing this gap — either
by sharpening the lower bound or by constructing a non-linear
encoder that matches the $2^{-2R/\deff}$ rate — is a concrete
open problem that we view as beyond the scope of this work.
Remark~\ref{rem:sharedv-gap} elaborates.

\paragraph{Scope of the hard distribution.} The lower bound is
established against $P^\star$ (Stiefel-random $V_t$, arbitrary
task-dependent spectra). Real LoRA fine-tunes exhibit structure
not captured by iid Stiefel random sampling: tasks that share
semantic similarity produce correlated subspaces; fine-tuning
induces anisotropic curvature; and the full-dimensional backbone
of the base model introduces rank-$d$ modes the bound ignores.
Extending the lower bound to structured task distributions, or
supplying a second construction that exploits task correlation
to beat our generic upper bound, are natural directions. Until
such extensions exist, the theorems of this paper should be read
as characterizing the \emph{information-theoretic} floor for
merging, not a tight prediction of any specific real-LLM merging
experiment.

\paragraph{What the real-LLM experiments add.} Section~\ref{sec:exp}
makes the scope caveat above concrete. We measure $\deff = Tr$ in
every layer of all four base models -- the task subspaces are
linearly independent, so the predicted floor is exactly zero --
yet real merged adapters incur $0.10$--$0.22$ nats/token of
worst-task excess. The gap is therefore not floor-driven; it is
the slack between current heuristics and the rate-distortion
optimum, and it is what a better merging algorithm could close.
The cross-model ordering of this slack is \emph{not} explained by
subspace overlap: a soft effective-dimension (participation ratio
of the stacked subspace) is uniform to within $0.7\%$ across
architectures that differ in merging difficulty by $2.2\times$,
and a Llama-family model (Yi-1.5-9B) merges more cleanly than
Llama-3.1 itself. This points to the loss curvature $H_t$ -- which
the bound takes as given -- and the pretraining/instruction-tuning
recipe that shapes it, as the operative quantities, and motivates
extending the theory to predict $H_t$ from training rather than
treating it as input.

\paragraph{Practical implications.} Two guidance points for
practitioners follow from the theorems. (i) The floor
$B^2(1 - \deff/(Tr))$ identifies when rate is wasted: if tasks
share substantial subspace overlap, no merging algorithm closes
the floor and additional bits buy little. This gives an
operational criterion for deciding when to merge at all versus
when to serve tasks from separate adapters. (ii) The
compression term $B^2 (\deff/(Tr)) 2^{-2R/\deff}$ is matched by
an explicit algorithm that runs in $O(\deff^2)$ time per merge
(the Gaussian-QR rotation dominates). For LoRA with $\deff = Tr$
in the thousands, this is negligible compared with the merging
overhead of existing heuristics; deployment cost is not a
blocker.

\paragraph{Known limitations.}
We list eight concrete limitations of the present results, in
descending order of how much they constrain the conclusions.
\begin{enumerate}[leftmargin=*,topsep=2pt,itemsep=2pt]
\item \textbf{Single LoRA rank.} All experiments use $r = 16$. The
log-$T$ slope $\rho \approx 1/\sqrt{r}$ predicted by the synthetic
sweep (\S\ref{subsec:exp-T-synthetic}) is consistent with the $r = 16$
measurements but not directly verified at other ranks; a focused
$r \in \{4, 8, 32, 64\}$ sweep is the natural extension.

\item \textbf{Four-base sweep.} Llama-3.1-8B-Instruct,
Mistral-7B-v0.3, Qwen-2.5-7B, and Yi-1.5-9B-Chat span an instructive
range of pretraining recipes and saturation regimes, but exclude the
Llama-3.3, Gemma-2, and Phi families. The cross-model story
(\S\ref{subsec:exp-T-scaling}, Figure~\ref{fig:hero}) is robust on the
two bases tested at $T \in \{2, 4, 7\}$; broader validation is
forthcoming.

\item \textbf{Seed coverage is uneven.} The $T = 4$ matrix
(\S\ref{subsec:exp-multiseed}) is run with two seeds plus bootstrap
CIs. The cross-model $T$-scaling sweep (\S\ref{subsec:exp-T-scaling})
uses $4$ random task subsets per $T$ but a single training seed per
adapter. The added baselines (\S\ref{subsec:exp-e10-e11}) and the
quadratic-bridge probe are seed-$1$ single cells.

\item \textbf{rd-encoder ridge tested on NLL only.} The downstream
metric verification of \S\ref{subsec:exp-downstream-gsm8k}
(GSM8K \textsc{em} and HumanEval pass@1) is run on the five matrix
methods, not on rd-encoder ridge. The R1 recommendation in
\S\ref{sec:practical-recommendations} therefore extrapolates from NLL
robustness, and we flag this gap explicitly.

\item \textbf{Single-cell outlier with no surviving mechanism.} The
(Llama-3.1-Instruct, GSM8K \textsc{em}) cell shows Spearman $\rho =
-0.60$ between NLL and accuracy across the five matrix methods, against
$\rho \geq +0.72$ on the other seven (base $\times$ metric) cells we
measured. Both pre-registered candidate mechanisms (chat-template
emission, cross-metric sign-election) were tested in
\S\ref{subsec:exp-downstream-gsm8k} and falsified. The outlier remains
empirically unexplained and is the one cell where we recommend running
$\geq 2$ downstream metrics in deployment.

\item \textbf{Quadratic bridge window is base-dependent.} E11
(\S\ref{subsec:exp-e11}) and E11b (Table~\ref{tab:e11b-llama-finer})
confirm rd-encoder ridge approaches the Fisher-quadratic optimum on
both bases tested, but the visible $\alpha$ window differs: Yi-Chat
shows the bridge in the small-$\alpha$ limit (ratio $\to 0.90$ at
$\alpha = 0.1$), Llama-Instruct in an intermediate window
$\alpha \in [0.15, 0.25]$ (ratio $0.83$ at $\alpha = 0.20$), bracketed
by ridge-floor over-saturation below and full-scale departure above.
The bridge holds on both, but the location of the window is
base-specific.

\item \textbf{Td2 sign-election threshold is order-of-magnitude.}
Appendix~\ref{app:sign-election-threshold} predicts the TIES win-share
range threshold $R^\star \in [0.025, 0.075]$ from a perturbation
analysis of the TIES update; the constants $c_{\mathrm{bias}}$,
$c_{\mathrm{sign}}$ are estimated by spectral analogy, not derived
from first-principles Fisher information. The two empirical points
(Yi $R = 0.077$, Llama $R = 0.028$) bracket the predicted range but
do not falsify it in the strict sense.

\item \textbf{Concurrent merge methods not exhaustively surveyed.}
The added baselines of \S\ref{subsec:exp-e10} (Fisher-magnitude,
DELLA) cover the two most commonly requested reviewer comparisons.
AdaMerging, RegMean, DARE-TIES and 2026-era methods such as
DO-Merging and 1-bit-Merging are discussed in
\S\ref{sec:related_work} but not all re-run on our matrix; the
empirical headline (rd-encoder ridge ties or beats TA and TIES on
every base at $T = 4$) is robust to the absence of any single
baseline, but a comprehensive 10+ method comparison is a fair
follow-up.
\end{enumerate}

\paragraph{Closing.} We have given LoRA merging the information-
theoretic reference point it had been missing. Worst-task NLL excess
splits cleanly into a task-overlap floor and a compression term that
decays exponentially in the bit budget; an explicit encoder closes
the gap up to a logarithmic constant in $T$. Empirically, real LoRA
task subspaces are linearly independent in every layer of every base
we tested, so the floor is zero --- everything current methods lose
is recoverable algorithmic slack rather than an information-theoretic
limit, and the regularized variant of our encoder shows how to claim
it. The practitioner-facing summary is the five-rule decision tree of
\S\ref{sec:practical-recommendations}; the most actionable single
result is the base-saturation diagnostic that predicts TIES inversion
from a one-minute CPU audit of any trained adapter cohort. We hope
both contribute to a phase of model-merging research that is judged
not only against the next benchmark, but against a theoretical
ceiling that can be measured.
